\documentclass{article}
\usepackage[ngerman]{babel}
\usepackage[utf8]{inputenc}
\usepackage[T1]{fontenc}
\usepackage{amsmath}
\usepackage{xurl}
\usepackage{hyperref}

\title{p-q-Formel und Differenzierbarkeit}
\author{Karl Kuhne}
\date{\today}

\begin{document}
\maketitle

\section{p-q-Formel}
Eine quadratische Gleichung in Normalform lautet
\[
  x^2 + px + q = 0.
\]
Ihre Lösungen berechnet man mit der p-q-Formel:
\[
  x_{1,2} = -\frac{p}{2} \pm \sqrt{\left(\frac{p}{2}\right)^2 - q}.
\]
Der Ausdruck unter der Wurzel entscheidet, ob es zwei, eine oder keine reelle Lösung gibt.

\subsection{Beispiel}
Für
\[
  x^2 - 4x + 3 = 0
\]
gilt $p=-4$ und $q=3$. Also:
\[
  x_{1,2}=2 \pm \sqrt{4-3}=2 \pm 1.
\]
Damit erhält man $x_1=3$ und $x_2=1$.

\section{Differenzierbarkeit}
Eine Funktion ist an einer Stelle differenzierbar, wenn dort der Grenzwert des Differenzenquotienten existiert:
\[
  f'(x)=\lim_{h\to 0}\frac{f(x+h)-f(x)}{h}.
\]
Die Ableitung beschreibt die Steigung der Tangente an den Graphen der Funktion.

\subsection{Zusammenhang mit Stetigkeit}
Jede differenzierbare Funktion ist an der betrachteten Stelle stetig. Umgekehrt muss eine stetige Funktion aber nicht differenzierbar sein.

\subsection{Beispiel}
Für $f(x)=x^2$ gilt:
\[
  f'(x)=2x.
\]
An der Stelle $x=3$ ist die Steigung also $f'(3)=6$.

\section*{Quellen}
\begin{itemize}
  \item \url{https://de.wikipedia.org/wiki/Quadratische_Gleichung#p-q-Formel}
  \item \url{https://de.wikipedia.org/wiki/Differentialrechnung#Differenzierbarkeit}
\end{itemize}

\end{document}